\section{Exponential control of cumulative collision counts}
\label{sec:cumulative-returns}
The avoidance estimate controls more than a single return. It also
controls the time of the $n$th visit to the section, without separating
the successive return blocks. This observation improves the truncation
budget in \eqref{eq:physical-count-truncation} and supplies exponential
moments on a fixed complex neighborhood for the entire physical record.

Retain $z_0,c,C$ from Proposition~\ref{prop:soft-killing}. Write
$N_{n,R}=\sum_{j<n}r_R^*\circ(F_R^*)^j$, and enlarge a fixed constant
$D\ge1$, if necessary, so that
\begin{equation}\label{eq:cumulative-record-domination}
 |J_{n,R}|_1\le D N_{n,R},\qquad Q_{n,R}\le D N_{n,R}.
\end{equation}
Such a constant exists because each physical flight has uniformly
bounded length and lattice displacement. All probabilities and
expectations in this section are with respect to $\nu_R^*$.

\begin{theorem}[The time of the $n$th actual return]
\label{thm:cumulative-return-tail}
There are constants $A\ge1$, $a,c>0$, independent of $n$, $L$ and $R$,
such that
\begin{equation}\label{eq:cumulative-return-tail}
 \nu_R^*\{N_{n,R}>L\}\le A\exp(an-cL)
       \qquad(n\ge1,\ L\ge0).
\end{equation}
One may take $a=z_0$ and the exponent $c$ of the soft-killing estimate.
Put $v=a/c$. For $0<s<c$ there is a constant $B_s<\infty$, uniform in
$n,R$, for which
\begin{equation}\label{eq:cumulative-exponential-moment}
 \int e^{s(N_{n,R}-vn)_+}\dd\nu_R^*\le B_s,
 \qquad \int e^{sN_{n,R}}\dd\nu_R^*\le B_s e^{svn}.
\end{equation}
Consequently, for every $0<\lambda<c/D$,
\begin{equation}\label{eq:fixed-strip-record-moment}
 \int e^{\lambda Q_{n,R}}\dd\nu_R^*
          \le B_{\lambda D}e^{\lambda Dvn}.
\end{equation}
The constants $B_s$ are locally bounded for $0\le s<c$, with $B_0=1$.
The number $v$ is an upper envelope, not the Kac mean $c_*^{-1}$.
\end{theorem}
\begin{proof}
For an integer $L\ge0$ let
\[
 H_{L,R}(x)=\sum_{j=1}^{L}\one_{Y_R^*}(T_R^j x).
\]
Since $N_{n,R}$ is the physical collision time of the $n$th return,
$N_{n,R}>L$ if and only if $H_{L,R}<n$ on $Y_R^*$. The fixed smooth
weight satisfies $0\le\chi\le\one_{Y_R^*}$. Thus, on this event,
$\sum_{j=1}^{L}\chi(T_R^j x)\le n-1$. In particular, pointwise on
$Y_R^*$,
\[
 \one_{\{N_{n,R}>L\}}
 \le e^{z_0(n-1)}
       \exp\left(-z_0\sum_{j=1}^{L}\chi(T_R^j x)\right).
\]
Integrate on $Y_R^*$, enlarge the domain to $M$, and use invariance to
shift the sum by one collision. Equation~\eqref{eq:soft-killing} gives
\[
 \nu_R^*\{N_{n,R}>L\}\le c_*^{-1}C e^{z_0(n-1)-cL}.
\]
For real $L$ replace it by $\lfloor L\rfloor$ and enlarge the constant
by $e^c$. This proves \eqref{eq:cumulative-return-tail}. It uses a single
orbit's visit count, not a product of probabilities of return blocks.

It follows that, for $u\ge0$,
$\nu_R^*\{N_{n,R}>vn+u\}\le A e^{-cu}$. For a nonnegative random
variable $X$, Tonelli's theorem gives
\[
 \int e^{sX}\dd\nu_R^*
       =1+s\int_0^\infty e^{su}\nu_R^*\{X>u\}\dd u.
\]
Apply this with $X=(N_{n,R}-vn)_+$ to obtain
$B_s=1+As/(c-s)$. Since $N_{n,R}\le vn+X$, the second bound in
\eqref{eq:cumulative-exponential-moment} follows. Finally apply
\eqref{eq:cumulative-record-domination} with $s=\lambda D$.
\end{proof}

We now apply the tail to the actual record measure. For a measurable
insertion $w_{n,R}$ define the finite signed or complex measure
\[
 \mu_{n,R}^{w}=(J_{n,R})_*(w_{n,R}\nu_R^*),\qquad
 \mu_{n,R}^{w,L}=(J_{n,R})_*
       (w_{n,R}\one_{\{N_{n,R}\le L\}}\nu_R^*).
\]
The transform uses the convention $\widehat\mu(\omega)=
\int e^{i\omega\cdot j}\dd\mu(j)$, with
$\omega=(u_1,u_2,s,b)\in\T^3\times\R$.

\begin{proposition}[Weighted finite-band truncation]
\label{prop:linear-count-budget}
Fix a nonnegative integer $d$ and suppose
$|w_{n,R}|\le M(1+Q_{n,R})^d$. For every multi-index $\gamma$ there is a
constant $C_{d,\gamma}$, independent of $n,R,L,w$, such that
\begin{align}\label{eq:weighted-count-truncation}
 \|\mu_{n,R}^{w}-\mu_{n,R}^{w,L}\|_{\TV}
  &\le C_{d,0}M(1+L)^d e^{an-cL},\\
 \sup_{\omega\in\T^3\times\R}
 \left|\partial_\omega^\gamma
   (\widehat\mu_{n,R}^{w}-\widehat\mu_{n,R}^{w,L})(\omega)\right|
  &\le C_{d,\gamma}M(1+L)^{d+|\gamma|}e^{an-cL}.\nonumber
\end{align}
For a bounded insertion ($d=0$), a frequency region of finite volume
$V>0$, and $\varepsilon>0$, the sufficient integer cutoff for integrated
transform error at most $\varepsilon$ is
\begin{equation}\label{eq:linear-count-budget}
 L\ge\left\lceil\frac{a}{c}n+
       \frac1c\max\left\{0,\log\frac{C_{0,0}MV}{\varepsilon}\right\}
        \right\rceil.
\end{equation}
If $V_n\le V_0e^{\kappa n}$ and the required error is $e^{-\delta n}$,
then a cutoff linear in $n$ suffices. This remains true for each fixed
$d,\gamma$, with any fixed positive slack in the exponential budget.
\end{proposition}
\begin{proof}
For an integer $m\ge0$, integration of the tail gives
\begin{align*}
 &\int (1+N_{n,R})^m\one_{\{N_{n,R}>L\}}\dd\nu_R^*\\
 &\quad=(1+L)^m\nu_R^*\{N_{n,R}>L\}
    +m\int_L^\infty(1+t)^{m-1}\nu_R^*\{N_{n,R}>t\}\dd t\\
 &\quad\le C_m(1+L)^m e^{an-cL}.
\end{align*}
The integral term is zero when $m=0$. For $m\ge1$, the last inequality
follows by writing $t=L+u$ and using
$1+L+u\le(1+L)(1+u)$. Pushforward does not increase total variation.
For the derivative estimate differentiate under the integral: its
absolute integrand is at most
$M(1+Q_{n,R})^d|J_{n,R}|_1^{|\gamma|}
\one_{\{N_{n,R}>L\}}$. Equation~\eqref{eq:cumulative-record-domination}
and the preceding tail integral prove both estimates.

Multiplication by $V$ bounds the integral of the transform error. Solving
$C_{0,0}MV e^{an-cL}\le\varepsilon$ gives
\eqref{eq:linear-count-budget}. More generally choose a slope
$\ell>(a+\kappa+\delta)/c$. Then, for $L_n=\lceil\ell n\rceil$,
\[
 V_n(1+L_n)^{d+|\gamma|}e^{an-cL_n}
       \le C(1+n)^{d+|\gamma|}e^{-(\delta+\epsilon_0)n}
\]
for some $\epsilon_0>0$. A fixed additive increase of the cutoffs absorbs
the remaining constant and polynomial, proving the last assertion.
\end{proof}

\begin{remark}
This improves the sufficient cutoff of order
$n\log(V/\varepsilon)$ obtained from
\eqref{eq:physical-count-truncation} to order
$n+\log(V/\varepsilon)$. On an exponentially growing frequency band the
former budget is quadratic in $n$, whereas the new one is linear.
Neither estimate bounds the number of critical words, the variation of
an inverse coarea Jacobian, or a pointwise raw density. In particular the
all-branch derivative sum in Lemma~\ref{lem:raw-residual-sum} remains a
separate requirement. The truncation removes only actual high-count
records in an error estimate; the limiting datum is still the full
unmodified physical record.
\end{remark}
